\section{Abstract}\label{sec:5-abstract}
Modern machine learning toolchains built on differentiable programming — PyTorch, JAX~\cite{bradbury_jax_2018} — have transformed computational physics by enabling gradient-based optimization over physical simulators.
In gravitational wave cosmology, the forward model mapping source parameters to the observed SGWB spectral density is a smooth, differentiable function of those parameters.
Yet existing PTA analysis frameworks (ENTERPRISE~\cite{ellis_enterprise_2020}, PTArcade~\cite{mitridate_ptarcade_2023a}) treat this model as a black box, using gradient-free samplers (MCMC) and discarding the information in the Jacobian of the forward model.
Simulation-Based Inference methods~\cite{cranmer_frontier_2020}, while powerful, also do not exploit the differentiability of the simulator.

This proposal develops a JAX-based~\cite{bradbury_jax_2018} differentiable forward simulator for SGWB signals in the PTA band, and investigates three inference strategies that exploit gradients: (1) gradient-based MCMC (NUTS/HMC~\cite{hoffman_nuts_2014}) via NumPyro~\cite{phan_composable_2019}; (2) variational inference (VI) using normalizing flows~\cite{papamakarios_normalizing_2021} with gradient-optimized parameters; and (3) Neural Posterior Estimation (NPE) with gradient-informed sequential simulation (active learning)~\cite{cranmer_frontier_2020}.
The result is both a software contribution (a publicly available differentiable SGWB simulator) and a methodological comparison — identifying which inference strategy is most simulation-efficient for cosmic string parameter estimation under realistic PTA noise~\cite{afzal_nanograv_2023}.

This proposal is the most technically ambitious and software-forward of the five.
It is appropriate if the student's primary interest is in the architecture of inference systems and differentiable scientific computing, with the GW physics as the application domain.


\section{Objectives}\label{sec:5-objectives}
\begin{itemize}
    \item Implement a fully differentiable SGWB forward model for cosmic strings~\cite{martins_extending_2002,blanco-pillado_number_2014} and a power-law SGWB in JAX~\cite{bradbury_jax_2018}, supporting \texttt{jit}, \texttt{grad}, \texttt{vmap}, and \texttt{pmap} transformations;
    \item Deploy gradient-based MCMC (HMC/NUTS~\cite{hoffman_nuts_2014} via NumPyro~\cite{phan_composable_2019}) on the differentiable model and benchmark convergence against gradient-free MCMC (PTMCMCSampler~\cite{ellis_enterprise_2020});
    \item Implement a normalizing flow posterior~\cite{papamakarios_normalizing_2021} (using \texttt{nflows} or \texttt{flowjax}) trained via maximum likelihood on simulated $(\theta, x)$ pairs, and compare to TMNRE~\cite{miller_truncated_2021};
    \item Implement a gradient-informed active learning loop for sequential simulation~\cite{cranmer_frontier_2020}, reducing the number of simulator calls needed to achieve target posterior accuracy;
    \item Release the differentiable simulator as a standalone JAX~\cite{bradbury_jax_2018} library compatible with the broader NumPyro~\cite{phan_composable_2019}/BlackJAX~\cite{cabezas_blackjax_2024} inference ecosystem.
\end{itemize}


\section{Methods}\label{sec:5-methods}
\begin{itemize}
    \item The VOS network evolution equations~\cite{martins_extending_2002} and the $\Omega_{\rm gw}(f)$ integral are implemented in JAX~\cite{bradbury_jax_2018} as pure functions, verified to be differentiable through \texttt{jax.grad} by comparison with finite-difference derivatives.
    \item Automatic differentiation through the ODE solver is achieved via the \texttt{diffrax} library~\cite{kidger_neural_2022} (Kidger 2021), which provides differentiable Dormand-Prince and Kvaerno solvers in JAX.
    \item For gradient-based MCMC, the PTA likelihood (Gaussian approximation to the free-spectrum posterior~\cite{johnson_nanograv_2024}) is implemented in NumPyro~\cite{phan_composable_2019} and NUTS~\cite{hoffman_nuts_2014} is run using BlackJAX~\cite{cabezas_blackjax_2024}, with wall-clock time and effective sample size compared to PTMCMCSampler.
    \item For normalizing flows~\cite{papamakarios_normalizing_2021}, a masked autoregressive flow is trained offline on $(\theta, \text{summary}(x))$ pairs, then used as an amortized posterior. Posterior quality is assessed via simulation-based calibration~\cite{talts_validating_2018}.
    \item The gradient-informed active learning loop selects new simulation points from regions of high posterior uncertainty (estimated via an ensemble of normalizing flows~\cite{papamakarios_normalizing_2021}), reducing simulation budget by targeted prior refinement~\cite{cranmer_frontier_2020}.
    \item All experiments are run on the group's GPU cluster nodes; the JAX~\cite{bradbury_jax_2018} implementation enables transparent multi-GPU distribution via \texttt{pmap}.
\end{itemize}